\documentclass[11pt]{article}

\usepackage[margin=1in]{geometry}
\usepackage{amsmath,amssymb,amsthm}
\usepackage{booktabs}
\usepackage{graphicx}
\usepackage{subcaption}
\usepackage{xcolor}
\usepackage{hyperref}
\usepackage{enumitem}
\usepackage{microtype}

\hypersetup{
    colorlinks=true,
    linkcolor=blue,
    citecolor=blue,
    urlcolor=blue
}

\newcommand{\KAN}{\mathrm{KAN}}
\newcommand{\MLP}{\mathrm{MLP}}
\newcommand{\VarF}{\mathrm{VarF1}}
\newcommand{\IntF}{\mathrm{IntF1}}
\newcommand{\MSE}{\mathrm{MSE}}
\newcommand{\supp}{\mathrm{supp}}
\newcommand{\E}{\mathbb{E}}

\title{
Diagnosing and Mitigating High-Dimensional Interaction Failures\\
in Kolmogorov--Arnold Network Explanations
}

\author{
Zicong Geng \\
University of Michigan \\
\texttt{zicong@umich.edu}
\and
Adrian Corrales \\
\texttt{[email]}
}

\date{\today}

\begin{document}
\maketitle

\begin{abstract}
Kolmogorov--Arnold Networks (KANs) are often presented as more interpretable alternatives to standard neural networks because their learned edge functions can be inspected directly.
However, it is unclear when these learned structures faithfully recover the variables and interactions of the underlying target function, especially in high-dimensional regimes.
We study this question through synthetic structure-recovery benchmarks where the ground-truth active variables and interaction terms are known.
We find that standard KAN explanations recover structure in clean low-dimensional settings, but fail on high-dimensional interaction functions: the model recovers main effects while missing interaction-only variables such as $x_2,x_3$ and the interaction $(x_2,x_3)$.
Oracle explanations computed directly on the true function recover the correct structure, showing that the evaluation metric is not the main source of failure.
Input-level and KAN path-level interventions further show that the trained high-dimensional KAN does not functionally rely on the true interaction variables, indicating a learned-representation failure rather than merely a post-hoc explanation failure.
Finally, support-controlled screening ablations show that feature screening rescues high-dimensional KAN recovery only when it retains the true interaction variables; random dimensionality reduction and interaction-excluding screens fail.
These results suggest that KAN explanation fidelity is highly sensitive to irrelevant-feature contamination and correlated proxy structure, and that faithful KAN explanations require explicit validation beyond visual inspection of learned edge functions.
\end{abstract}

\section{Introduction}

Kolmogorov--Arnold Networks (KANs) replace fixed scalar weights with learned univariate functions on edges.
This architectural choice suggests a natural path toward interpretability: if an edge function or path is active, it may correspond to a meaningful contribution of the associated input variable.
However, the fact that a network is visually inspectable does not guarantee that the induced explanation is faithful to the ground-truth data-generating mechanism.

This paper studies a basic but important question:

\begin{quote}
\emph{When do KAN explanations faithfully recover the variables and interactions of the target function?}
\end{quote}

We focus on controlled synthetic settings where the true formula is known.
This lets us distinguish several different sources of explanation failure:
\begin{enumerate}[leftmargin=2em]
    \item the explanation metric itself may fail to identify the true structure;
    \item the post-hoc explanation of the trained KAN may fail, even if the model learned the correct functional dependence;
    \item the trained KAN may not have learned the true functional dependence at all;
    \item correlated proxy variables may create predictively useful but formula-unfaithful explanations.
\end{enumerate}

Our main finding is that high-dimensional KAN explanation failure is not merely a post-hoc explanation artifact.
In high-dimensional sparse interaction settings, trained KANs often do not functionally rely on the true interaction variables.
This is confirmed both by input-level interventions and by KAN-specific path interventions.
At the same time, the failure is not an unavoidable expressivity limitation: if a supervised screening stage preserves the true interaction variables, KAN recovery can be restored.

\paragraph{Contributions.}
This paper makes four empirical and methodological contributions.
\begin{enumerate}[leftmargin=2em]
    \item We introduce a structure-recovery benchmark for KAN explanations, measuring both active-variable recovery and interaction recovery.
    \item We propose a three-layer failure-attribution framework: oracle explanations on the true function, explanations of the trained KAN, and model/path-level interventions.
    \item We show that high-dimensional KANs can fail to form functional paths for true interaction variables, even when oracle explanations recover the correct structure.
    \item We show through support-controlled screening ablations that feature screening rescues high-dimensional KAN recovery only when the true interaction variables are retained.
\end{enumerate}

\section{Problem Setup}

\subsection{Ground-Truth Structure Recovery}

We consider supervised regression data
\[
    y = f(x) + \varepsilon, \qquad x \in \mathbb{R}^d,
\]
where the formula $f$ is known in synthetic experiments.
Each function has a known set of active variables
\[
    S^\star \subseteq \{0,\dots,d-1\},
\]
and a known set of interacting variable pairs
\[
    I^\star \subseteq \{(i,j):0 \leq i < j < d\}.
\]

The main benchmark function used throughout the paper is
\[
    f_{\mathrm{core}}(x)
    =
    \sin(2\pi x_0) + x_1^2 + c x_2 x_3,
\]
where $c=1$ unless otherwise stated.
The true active set is
\[
    S^\star = \{0,1,2,3\},
\]
and the true interaction set is
\[
    I^\star = \{(2,3)\}.
\]

We evaluate three regimes:
\begin{itemize}[leftmargin=2em]
    \item \textbf{Clean low-dimensional}: $d=20$, $c=1$.
    \item \textbf{High-dimensional core}: $d=100$, $c=1$.
    \item \textbf{Strong interaction}: $d=100$, $c=5$.
    \item \textbf{Correlated proxy}: $d=100$, with proxy variables such as $x_4 \approx x_0$.
\end{itemize}

\subsection{Explanation Metrics}

For a trained model $\hat f$, we compute variable-level explanation scores using gradient-based and permutation-based importance.
Given a score vector $a \in \mathbb{R}^d$, we select the top-$|S^\star|$ variables and compute precision, recall, and F1 against $S^\star$.
We also report AUROC and AUPRC when useful.

For interaction recovery, we use Hessian-based pair scores:
\[
    h_{ij}
    =
    \E_x \left| \frac{\partial^2 \hat f(x)}{\partial x_i \partial x_j} \right|.
\]
We select the top-$|I^\star|$ pairs and compute interaction F1 against $I^\star$.

\subsection{Failure Attribution Framework}

We evaluate three levels of explanation fidelity.

\paragraph{Oracle explanation.}
We compute explanation scores directly from the true function $f$.
If oracle explanations recover $S^\star$ and $I^\star$, then the metric is capable of detecting the true structure.

\paragraph{Trained-KAN explanation.}
We train KANs and compute the same explanation metrics on $\hat f_{\KAN}$.
This measures whether the trained model explanation recovers the known structure.

\paragraph{Model and path interventions.}
We intervene on the trained model to test functional reliance.
Input-level interventions permute or zero selected variables.
KAN-specific path interventions delete all input-to-hidden KAN edges from a feature, or delete feature-pair paths through shared hidden units.
If deleting $x_i$ paths does not change test MSE, then the trained KAN does not functionally rely on $x_i$.

\section{Experimental Design}

\subsection{KAN Training}

Unless otherwise stated, we train KANs using the same default architecture and training procedure across settings.
The baseline/default KAN uses width $[d,8,1]$, spline grid size $5$, spline order $3$, and LBFGS optimization.
We also consider tuned KAN hyperparameters when available.

\subsection{Screening and Support-Controlled Ablations}

To test whether high-dimensional failure is caused by irrelevant-feature contamination, we train KANs after reducing the input dimension to $M$ selected variables.
We compare the following screening modes:
\begin{itemize}[leftmargin=2em]
    \item \textbf{Random}: choose $M$ variables uniformly at random.
    \item \textbf{Oracle support}: force inclusion of all true variables $S^\star$, then fill the remaining slots randomly.
    \item \textbf{Exclude interaction}: choose $M$ variables while excluding $x_2,x_3$.
    \item \textbf{RF}: select the top-$M$ variables by random-forest feature importance.
    \item \textbf{RF exclude interaction}: rank variables by random forest but explicitly exclude $x_2,x_3$.
\end{itemize}

This design separates two hypotheses:
\[
    \text{screening helps because it reduces dimension}
\]
versus
\[
    \text{screening helps because it retains the true interaction variables.}
\]

\section{Results}

\subsection{Result 1: Raw KANs Fail on High-Dimensional Interaction Recovery}

In clean low-dimensional settings, KANs recover both active variables and the true interaction.
However, in $d=100$, raw KANs recover main-effect variables such as $x_0,x_1$ much more reliably than interaction-only variables $x_2,x_3$.
The interaction $(2,3)$ is often missed.

\paragraph{Interpretation.}
This suggests that KAN explanations are not uniformly faithful even in synthetic settings where the target formula is simple.
The failure is especially visible for interaction-only variables whose marginal signal is weak.

\begin{figure}[t]
    \centering
    \includegraphics[width=0.85\linewidth]{results/master_figures/master_variable_recovery.png}
    \caption{
    Variable recovery for raw KAN, stability-selected KAN, and RF-screened KAN.
    Raw KANs fail in high-dimensional settings, while RF-screened KANs recover variables when the screen retains the true support.
    }
    \label{fig:variable-recovery}
\end{figure}

\begin{figure}[t]
    \centering
    \includegraphics[width=0.85\linewidth]{results/master_figures/master_interaction_recovery.png}
    \caption{
    Interaction recovery for raw KAN, stability-selected KAN, and RF-screened KAN.
    Stability selection alone cannot recover interactions that the trained KAN failed to learn, while RF screening can rescue recovery when it preserves $x_2,x_3$.
    }
    \label{fig:interaction-recovery}
\end{figure}

\subsection{Result 2: The Failure Is Model-Level, Not Merely Explanation-Level}

Oracle explanations computed directly on the true function recover the correct variables and interactions even in high dimension.
Thus, the explanation metric is not the main bottleneck.

We then perform input-level and KAN path-level interventions.
In clean $d=20$ settings, deleting paths from $x_2,x_3$ significantly increases test MSE.
In high-dimensional $d=100$ settings, deleting $x_2,x_3$ paths has little effect, while deleting $x_0$ paths has a large effect.

\paragraph{Interpretation.}
The trained high-dimensional KAN does not merely hide the interaction from the explanation method.
It often fails to form a functional internal path for the true interaction variables.

\begin{figure}[t]
    \centering
    \includegraphics[width=0.85\linewidth]{results/master_figures/master_core_feature_path_reliance.png}
    \caption{
    Feature-path deletion in core $d=100$.
    Raw KAN relies strongly on $x_0$ but weakly on $x_2,x_3$.
    RF-screened KAN restores functional reliance on the true interaction variables.
    }
    \label{fig:core-feature-path}
\end{figure}

\begin{figure}[t]
    \centering
    \includegraphics[width=0.75\linewidth]{results/master_figures/master_core_pair_path_reliance.png}
    \caption{
    Pair-path deletion in core $d=100$.
    Raw KAN has weak reliance on the true interaction pair $(2,3)$.
    RF-screening restores pair-path reliance.
    }
    \label{fig:core-pair-path}
\end{figure}

\subsection{Result 3: Stability Selection Helps Variable Recovery but Cannot Restore Missing Interactions}

We next aggregate explanations across random seeds using stability selection.
For variable recovery, stability selection reduces random false positives and often improves F1.
However, interaction recovery remains poor when the trained KAN never learned the interaction.

\paragraph{Interpretation.}
Stability selection can reduce seed-level noise, but it cannot recover a structure absent from the trained model.
This supports the intervention result: missing interactions are a learned-representation failure, not merely an unstable explanation ranking.

\subsection{Result 4: Screening Rescues KAN Only When It Retains the True Interaction Variables}

Support-controlled screening ablations show that screening helps because it preserves the true interaction variables, not merely because it reduces input dimension.
Random screens fail.
Interaction-excluding screens fail.
Oracle-support screens recover structure.
RF screens recover structure when they retain $x_2,x_3$.
RF screens that explicitly exclude $x_2,x_3$ fail despite using the same supervised ranking procedure.

\paragraph{Interpretation.}
The key causal factor is support retention.
Irrelevant-feature contamination prevents KAN from learning faithful high-dimensional interaction structure, but the failure can be mitigated when screening removes irrelevant variables while retaining the true interaction endpoints.

\begin{figure}[t]
    \centering
    \includegraphics[width=0.9\linewidth]{results/support_controlled/clean_summary/support_controlled_interaction_f1_clean.png}
    \caption{
    Support-controlled ablation.
    Interaction recovery is high for oracle-support and RF screens, but low for random screens and interaction-excluding screens.
    This shows that screening helps because it retains the true interaction variables, not merely because it reduces dimensionality.
    }
    \label{fig:support-controlled}
\end{figure}

\begin{figure}[t]
    \centering
    \includegraphics[width=\linewidth]{results/support_controlled/clean_summary/support_controlled_interaction_f1_by_m.png}
    \caption{
    Interaction recovery across screening budget $M$.
    Core settings are rescued by support-preserving screens.
    The correlated-proxy setting remains more difficult, indicating an additional formula-fidelity challenge.
    }
    \label{fig:support-controlled-by-m}
\end{figure}

\subsection{Result 5: Correlated Proxies Remain a Separate Formula-Fidelity Challenge}

In correlated proxy settings, RF screening can restore some reliance on $x_2,x_3$, but proxy variables such as $x_4 \approx x_0$ remain functionally important.
Thus, the model may be predictively faithful but not formula-faithful.

\paragraph{Interpretation.}
This distinguishes predictive fidelity from formula-level fidelity.
A KAN explanation may correctly describe what the trained model uses, while still failing to recover the original symbolic mechanism.

\begin{figure}[t]
    \centering
    \includegraphics[width=0.85\linewidth]{results/master_figures/master_proxy_feature_path_reliance.png}
    \caption{
    Feature-path reliance in the correlated-proxy setting.
    Screening can restore true interaction-variable paths, but proxy-related paths remain important.
    }
    \label{fig:proxy-feature-path}
\end{figure}

\begin{figure}[t]
    \centering
    \includegraphics[width=0.75\linewidth]{results/master_figures/master_proxy_pair_path_reliance.png}
    \caption{
    Pair-path reliance in the correlated-proxy setting.
    The true pair $(2,3)$ can be partially restored, but proxy pair $(0,4)$ remains important.
    }
    \label{fig:proxy-pair-path}
\end{figure}

\section{Discussion}

\subsection{Prediction Accuracy Is Not Sufficient for Explanation Fidelity}

A low test MSE does not guarantee that a KAN explanation recovers the true generating structure.
A model may achieve good prediction by relying on main effects, proxies, or alternative representations, while missing interaction-only variables.
This motivates evaluating both predictive performance and structure recovery.

\subsection{KAN Interpretability Requires Intervention Validation}

Visual inspection of learned edge functions is insufficient.
A learned path should be considered explanatory only if deleting or intervening on that path changes the model's prediction.
KAN path interventions provide an architecture-specific way to test whether the claimed explanation is functionally used.

\subsection{Screening as Mitigation}

Feature screening is effective when it removes irrelevant dimensions while preserving true interaction endpoints.
However, screening can also fail if interaction-only variables have weak marginal signal, or if correlated proxies dominate the screening score.
Thus, screening is a mitigation strategy, not a full solution.

\section{Related Work}

\paragraph{Kolmogorov--Arnold Networks.}
KANs propose learnable univariate edge functions as an interpretable alternative to standard neural networks.
Our work studies when this interpretability claim holds under ground-truth structure recovery.

\paragraph{Explainable AI and faithfulness.}
Our results connect to broader concerns in XAI: explanations can be visually plausible or predictively useful while failing to match the causal or formula-level data-generating process.

\paragraph{Interaction detection.}
Our Hessian-based interaction recovery is related to methods for detecting non-additive dependencies.
The key distinction is that we evaluate whether a trained KAN actually learns and uses such interactions.

\paragraph{Stability selection and feature screening.}
Stability selection reduces seed-level noise, while feature screening can remove irrelevant-feature contamination.
Our support-controlled ablations show that support retention, not dimensionality reduction alone, explains screened KAN recovery.

\section{Limitations}

This paper currently focuses on synthetic functions where the true formula is known.
This is necessary for precise structure-recovery evaluation, but future work should test whether the same failure modes appear in real scientific or tabular datasets.
Our screening methods are supervised and may not be appropriate in all settings.
The correlated-proxy experiments show that even successful screening can preserve predictive proxies that obscure formula-level fidelity.
Finally, results may depend on KAN implementation details, optimization settings, and hyperparameter tuning.

\section{Conclusion}

We present a failure-attribution study of KAN explanation fidelity.
In high-dimensional sparse interaction settings, raw KANs can fail to recover interaction-only variables and the true interaction pair.
Oracle explanations show that the metric itself can recover the structure, while input and path interventions show that the trained KAN often does not functionally rely on the true interaction variables.
Stability selection reduces random explanation noise but cannot recover missing interactions.
Feature screening can rescue high-dimensional KAN recovery, but support-controlled ablations show that this works only when the true interaction variables are retained.
Correlated proxy settings remain a separate formula-fidelity challenge.
Together, these results suggest that KAN explanations require functional validation and careful treatment of high-dimensional feature contamination.

\bibliographystyle{plain}
\bibliography{references}

\appendix

\section{Additional Experimental Details}

\subsection{Synthetic Functions}

\paragraph{Core interaction.}
\[
    f(x)=\sin(2\pi x_0)+x_1^2+c x_2x_3.
\]

\paragraph{Correlated proxy.}
The target function uses $x_0,x_1,x_2,x_3$, while additional variables such as $x_4$ are correlated proxies for true variables.

\subsection{Evaluation Metrics}

Variable F1 is computed by selecting the top-$|S^\star|$ variables by explanation score.
Interaction F1 is computed by selecting the top-$|I^\star|$ pairs by Hessian interaction score.

\subsection{Path Intervention}

For a KAN with architecture $[d,h,1]$, the first layer contains input-to-hidden edge functions and the second layer contains hidden-to-output edge functions.
Feature-path deletion removes all first-layer edges from a selected input variable.
Pair-path deletion removes all first-layer edges from both variables in a pair.
Shared-hidden pair-edge deletion removes the two input-to-hidden edges from a pair into the same hidden unit.

\end{document}
